\section{Результаты имитационного моделирования}

\subsection{Исследование зависимости амплитуды побочных составляющих от разрядности ЦГ}

В данном пункте исследовалась зависимость амплитуды побочных (не связанных с полезным входным сигналом) спектральных составляющих, возникающих в спектре сигнала на выходе цифрового фильтра, от разрядности цифрового гетеродина ($n_g$). Для анализа данного явления параметр типа вычислений ЦГ был задан как \texttt{nco\_type = 'fixed'}. Разрядность гетеродина варьировалось в диапазоне от 2 до $n_a + 2$ бит, где $n_a = 9$ — количество бит АЦП (Мой номер по журналу является 4).

Оценка уровня побочных составляющих проводилась путем измерения динамического диапазона, свободного от паразитных составляющих (SFDR). Итоговый график зависимости представлен на рисунке \ref{fig:sfdr_vs_bits}.

\begin{figure}[H]
    \centering
    \includegraphics[width=0.8\textwidth]{sfdr 3 pynkt.png}
    \caption{Зависимость уровня подавления побочных составляющих (SFDR) от количества бит ЦГ}
    \label{fig:sfdr_vs_bits}
\end{figure}

\textbf{Выводы по результатам моделирования:}
Представленный график демонстрирует, что значение SFDR стабилизируется на уровне $\approx 70.7$ дБн, а все его изменения лежат в крайне узком диапазоне (от $70.62$ до $70.73$ дБн).

Математическая причина возникновения мощных побочных излучений в моделях ЦГ связана со спектральным наложением гармоник, порождаемых эффектами квантования амплитуды и усечения фазы. При экстремально низкой разрядности ($n_g=2$ бита) форма сигнала гетеродина максимально искажена, и амплитуда высших гармоник наибольшая — это отражается в локальном минимуме SFDR на графике.

Однако дальнейший рост SFDR жестко ограничивается шумовыми характеристиками предшествующего тракта. Базовый уровень аддитивного белого гауссова шума, заданный во входном сигнале (SNR = 70 дБ), а также шумы квантования 9-битного АЦП формируют плотную шумовую полку. Начиная с $n_g=3$ бит, паразитные гармоники гетеродина проваливаются под эту полку и перестают быть доминирующим источником помех. 

Незначительные колебания графика на уровне сотых долей децибела при $n_g \ge 3$ обусловлены статистическими флуктуациями спектральных отсчетов случайного шума при расчете преобразования Фурье, а не реальной динамикой побочных составляющих ЦГ. Таким образом, в условиях заданного шумового фона тракта повышение разрядности гетеродина свыше 3-4 бит не приводит к выигрышу в качестве демодулированного сигнала.